\section{Introduction}
\label{sec:intro}

Location-based social networks, such as Gowalla~\cite{cho2011gowalla}, let people share the places that they visit as check-ins, which together record how people move through a city. Individual mobility is highly predictable in principle~\cite{song2010limits}, so a service that anticipates the next move can prepare ahead of time, caching content where a user is heading~\cite{bastug2014edge} or provisioning capacity at the destination before demand arrives. Handover predictions already let cellular services adapt in advance at the network level~\cite{vielhaus2022handover}; at the city scale, check-in mobility analysis is likewise a design input for mobility-aware services~\cite{moura2025mobilityaware}.

Predicting the next place exactly, a single point of interest (POI), is hard and often more than a service needs. Two coarser questions are usually enough: the next category, the type of place such as food or shopping, and the next region, the part of the city that a user is heading to. We study whether one model should learn both at once.

Sharing one representation across tasks has a cost: in multi-task learning (MTL), the shared parameters can converge to a compromise optimal for neither task, helping one while hurting the other~\cite{caruana1997multitask}. Our earlier
work reported no consistent multi-task advantage for the paired category tasks and attributed it, in part, to this effect~\cite{silva2025mtlnet}; the useful question is where sharing helps and what it costs.

We propose two enhancements to that earlier model. First, we build a check-in-level representation: instead of one fixed vector per place, each check-in gets its own vector that holds its context (the time, nearby places, and recent visits).
This adds a fourth level, the check-in, beneath the place, region, and city levels of hierarchical graph infomax~\cite{huang2023hgi,velickovic2019dgi} (Fig.~\ref{fig:dataflow}). Second, we train one model that
predicts the next category and the next region in a single forward pass, with a shared trunk (a cross-attention stack where the two tasks exchange
semantic context) and a private spatial path for the region task (Fig.~\ref{fig:model}). We evaluate on two settings chosen to differ: five U.S. states of different sizes (Gowalla) and one non-U.S. city (Istanbul).\footnote{\label{fn:code}Code (model, representation, baselines, statistical tests): \url{https://github.com/VitorHugoOli/PoiMtlNet/tree/mobiwac}. Both data sources are public: the category-annotated Gowalla dump (\url{https://figshare.com/articles/dataset/gowalla_data/22126586}) and Massive-STEPS~\cite{wongso2025massivesteps}.}

Our contributions are as follows.
\begin{itemize}
  \item \textbf{A check-in-level representation.} Under a controlled comparison in which only the input
    changes, at one random initialization, it improves next-category prediction at every one of
    the six datasets and in every one of their folds, by $+0.23$ to $+6.29$ points of macro-F1,
    under one point at the three largest datasets, and a paired test separates the two
    representations at five of the six (Table~\ref{tab:substrate}). The complete system is also above every external baseline of Table~\ref{tab:results},
    on both tasks and at every dataset.
  \item \textbf{A single model for both tasks.} One model predicts the next category and the next
    region in one forward pass, reading both answers from a single saved model, so there is one
    artifact to train, version, and deploy. To our knowledge, fine-grained region as an end target
    of equal standing is underexplored (Section~\ref{sec:related-tasks}). The joint model outperforms a dedicated single-task region model at the
    two datasets with the largest region vocabularies and stays within a registered two-point
    margin at the other four (statistical non-inferiority, assessed with two one-sided tests). The joint model is also the larger model, so we do not attribute those gains to sharing
    (Section~\ref{sec:discussion}). On category it outperforms the dedicated model
    at one dataset, and the five remaining differences are equivalent to zero within half a point,
    even though the dedicated category model receives a per-dataset search
    (Table~\ref{tab:results}). The pairing of the two region gains with the two largest region vocabularies is an
    observation and not a law (Section~\ref{sec:results-part2}).
\end{itemize}

